% Page 20 — Alban en sa « légende dorée »
% À intégrer dans le document principal (XeLaTeX).
% Nécessite le package lettrine.

\begin{center}
    {\Large\itshape Alban en sa « légende dorée »}
\end{center}

\medskip

\lettrine[lines=2]{A}{lors}, comme le prodigue de la parabole, le malheur touche son âme,
les larmes jaillissent de ses yeux, et du cœur à ses lèvres monte la prière
à la Madone, cette oraison qu’on apprend tout petit sur les genoux d’une
pieuse mère :

\begin{center}
    O Notre Dame d’Alban -- Mère très aimée\\
    Tenez-nous toujours blottis contre votre cœur immaculé\\
    Afin que nous croissions, de jour en jour, en sagesse\\
    Et dans la crainte du Seigneur.
\end{center}

\medskip

\lettrine[lines=2]{A}{lain} n’est plus Alain. Alain est converti et le voici au soir d’une
terrible journée, à la nuit noire, demandant asile et protection au prieuré
bénédictin de Notre Dame d’Alban. Comment les disciples de Saint Benoît
auraient-ils pu se récuser devant les larmes et les prières d’un pareil
larron repentant ? Comment refuser pareil novice ? On tenta l’aventure
charitablement. Frère Alain prit donc le froc monacal. En attendant une
conversion totale et bien affermie, il rendrait à ses frères quelques menus
services, balayant la chapelle, désherbant le cimetière, arrosant le potager.
On verrait par la suite. Et la suite dépassa les espoirs les plus optimistes.
Qui eût pu, sous cette bure pacifique, reconnaître le terrible sire de
Nougaret ? Certainement pas les enfants…
